%% kapa: exam=1 round=29 year=2018 period=1 form=A subject=civil_law subject_ko=민법 seq=7 number=7
%% source: Q-Net 감정평가사 기출문제 (cst00309) artlSeq=5202808 file=extracted/1008/5202808/2018년 제29회 감정평가사 1차시험 기출문제/2018년 제29회 감정평가사 1차시험 기출문제/1교시 A형.pdf
%% conversion: kordoc 4.24.51 -> tools/gen1.py (auto); review=proofread-sonnet; flags=none
\begin{question}{7}
과실을 수취할 수 있는 자를 모두 고른 것은?
\begin{framedtext}
ㄱ. 질물의 과실에 대한 질권자 \par
ㄴ. 유치물의 과실에 대한 유치권자 \par
ㄷ. 점유물의 과실에 대한 선의의 점유자 \par
ㄹ. 토지전세권에서 토지의 과실에 대한 전세권설정자 \par
\end{framedtext}
\begin{choices}
  \item ㄱ, ㄴ
  \item ㄷ, ㄹ
  \item ㄱ, ㄴ, ㄷ
  \item ㄱ, ㄷ, ㄹ
  \item ㄴ, ㄷ, ㄹ
\end{choices}
\end{question}
\answer{3}
